\section{A three-cell laboratory: the hub participates even when it does not change}
The two-cell examples isolate curvature cleanly. A hub introduces another
question: can two actions interact even when the hub's final stock equals
its starting stock? Use
\begin{equation}
 z=(18,10,2),\quad x^*=(10,10,10),\quad\sigma=(2,2,2).
\end{equation}
Initial potential is sixteen and the marginals are $(2,0,-2)$.
Action A transfers four from Ridge to Hub; action B transfers four from
Hub to Vale. Their increments are
\begin{equation}
 \delta_A=(-4,4,0),\qquad\delta_B=(0,-4,4).
\end{equation}
Alone, A ends at $(14,14,2)$ and B alone ends at $(18,6,6)$.
Each endpoint has potential twelve, so each singleton value is four.
Together their increment is $(-4,0,4)$; their endpoint is $(14,10,6)$,
with potential four. The group value is twelve, not eight.

Calculate the cross term explicitly:
\begin{equation}
 \delta_A^TH\delta_B
 =\tfrac14[(-4)0+4(-4)+0(4)]=-4.
\end{equation}
The correction to the singleton sum is therefore plus four. It arises
because the hypothetical singleton endpoints separately disturb the hub
in opposite directions, whereas the joint endpoint does not. No
nonseparable potential was introduced. Shared-coordinate action structure
is sufficient to create the interaction.

\section{Three descriptions of the same aggregate endpoint}
Serialize A first, then B, keeping these quantities fixed and assuming
the simple stock constraints allow both. A contributes four and B then
contributes eight. Reverse the order: B contributes four and A then eight.
Both serialized field totals are twelve. Hence the positive correction
against hypothetical same-baseline singletons is not a positive advantage
against either of these actual sequential endpoints.

For the constant-rate common group path, the hub stays at ten throughout.
The gradient averaged along the path is $(3/2,0,-3/2)$. Therefore
\begin{equation}
 R_A=-(3/2,0,-3/2)\cdot(-4,4,0)=6,
 \qquad R_B=6.
\end{equation}
The split $(6,6)$ closes on twelve, just as the two sequential splits
$(4,8)$ and $(8,4)$ do. They refer to different descriptions of how
children contribute along a declared path. The equality of their totals
does not identify one split as causal, fair or institutionally required.

The reader should be able to explain all three descriptions without saying
that the same physical experiment has three contradictory outcomes. Only
one physical chronology or path is actually declared in any given model.
The other calculations are named comparisons or conventions.

\section{A shared-source laboratory with three children}
Now use four cells, all with reference ten and scale two, at
\begin{equation}
 z=(18,6,6,10).
\end{equation}
The total is forty and the initial potential is twelve. Three children
each withdraw two from the first cell, delivering respectively to cells
two, three and four. The final group state is $(12,8,8,12)$ and its
potential is two. Thus the group field value is ten.

The singleton values are five, five and three. All three share the source;
each pair has $\delta_a^TH\delta_b=1$. The singleton sum thirteen
therefore needs three corrections of minus one, yielding ten. A program
that only checks whether destination supports are disjoint will miss the
common source and return the wrong total.

For the common path, $\mu(z)=(2,-1,-1,0)$ and
$H\delta_G=(-3/2,1/2,1/2,1/2)$. The path-average gradient is
\begin{equation}
 \bar\mu=(5/4,-3/4,-3/4,1/4).
\end{equation}
Dotting it negatively with each child gives $(4,4,2)$. The values add
to ten. Notice that the child ending at an already-at-reference destination
still has a positive net value here: its source relief outweighs the
destination's added deviation. The account describes the combined physical
change, not only a destination label.

\section{Write the entire subset table before interpreting it}
For those same fixed increments, evaluate every subset at the same frozen
baseline:
\begin{center}
\begin{tabular}{cr}
\toprule Subset & Field value\\\midrule
$\varnothing$ & 0\\
A & 5\\
B & 5\\
C & 3\\
AB & 9\\
AC & 7\\
BC & 7\\
ABC & 10\\\bottomrule
\end{tabular}
\end{center}
Each pair coefficient is minus one. The third-order coefficient is
$10-9-7-7+5+5+3=0$. This cancellation illustrates the quadratic
membership structure. It is not a measured absence of all three-way
physical effects in nature.

Before calling the table a physical subset experiment, check that each
subset is admissible under the same protocol. If a coupled machine can
perform only all three children together, the singleton rows remain
algebraic counterfactuals, not admissible interventions. If a resolver
changes quantities separately for each subset, the table must be recomputed
with those quantities, and the fixed-increment degree argument may fail.
Do not fill missing physical rows with convenient numbers merely to make
an inversion formula available.

\section{An independence check with disjoint actions}
Let four cells start at $(14,6,14,6)$, again with reference ten and
scale two. One action sends two from cell one to cell two; another sends
two from cell three to cell four. Initial potential is eight. Each action
reduces its pair's potential from four to one, so each field value is
three. The group endpoint $(12,8,12,8)$ has potential two, giving six.

The increments have disjoint support, and the diagonal Hessian gives zero
cross term. Here independent singleton addition is valid. The lesson is
not that singleton addition is always wrong; it is that its validity has
a precise condition that can be checked. If both actions share a physical
machine limit, valuation may still add while joint permission fails.
Independence of field terms is not independence of every constraint.

\section{A positive group can have an unaffordable child owner}
Return to three cells, at $(14,6,10)$ with the same references and
scales. Let A send two from cell one to cell two and B send two from
cell two to cell three. The group endpoint is $(12,6,12)$.
Potential changes from four to three, so the group value is one.

The initial gradient is $(1,-1,0)$ and the net increment is
$(-2,0,2)$. The average gradient along the group path is
$(3/4,-1,1/4)$. Hence
\begin{equation}
 R_A=7/2,\qquad R_B=-5/2,\qquad R_A+R_B=1.
\end{equation}
This is an exact example of a positive group total with a negative
child attribution. If two different owners start with zero balances,
the proposed owner-wise rule refuses B's debit. If one owner holds both
children, same-event netting gives plus one and would pass that account
condition. Physical feasibility is the same in both descriptions.

The example is not an argument for selecting whichever ownership map
admits the group. It proves that the ownership map is a load-bearing model
choice. Scientific comparison requires that choice to be declared before
outcomes are inspected. It cannot be inferred from the favorable sign
of the group total.

\section{A missed factor can mimic an attractive improvement}
Suppose the two-cell lossless example is expanded to represent a third
physical resource consumed by pumping. The expanded potential might
contain a factor of that resource, and the accepted action might change
it. A quote that still uses only the water endpoints would omit a
changed term. Even if its arithmetic for water is perfect, it would no
longer be the complete exact local quote for the expanded declaration.

There are two legitimate ways to keep a teaching model simple. One is
to leave pumping energy outside its scope and state the limitation.
The other is to represent it and include the corresponding affected
factors or a separately justified residual burden. What is not legitimate
is to claim complete multi-carrier valuation while counting the same
physical effect in both places, or in neither place, without disclosure.

This is also why a fuel-to-electricity-to-water chain should remain a
linked sequence of atomic physical transformations when it is in scope.
The water endpoint does not tell the reader which upstream transformations
occurred or which loss boundary they crossed. A convenient aggregate
action label cannot repair missing physical representation.

\section{Recalibration is not physical work}
Take the original two-cell state $(18,2)$ and leave it unchanged.
With reference $(10,10)$ and scales two, potential is sixteen.
If both scales are revised to four, potential becomes four. The twelve-unit
difference is not an actor's physical improvement. No resource moved.
It is a change in the declared comparison geometry.

The active first programme freezes scales and references, so this event
is outside its action loop. A later adaptive-calibration model would need
an explicit parameter-change account. Otherwise an actor could appear
to earn value by changing the ruler rather than the world. Version
commitments protect the distinction between these operations; numerical
precision alone cannot.

\section{Build a receipt that another reader can reconstruct}
For the three-child source example, a minimal mathematical audit records
the ordered state $(18,6,6,10)$, the ordered reference and scales,
the three named increments, their sum, the before and after potentials,
the group value ten, and the chosen path values $(4,4,2)$. It also
records that the quantities are source-side quantities, not rates.

An implementation-grade record needs more: state and parameter identities,
the accepted physical constraints, path semantics, numerical policy,
and evidence that the executed increment matched the valued increment.
A capacity record additionally needs a declared owner map, prior balances
and a separately authorized update rule. Each layer answers a question
that the preceding arithmetic cannot answer by itself.

If a reader cannot determine whether the state list is ordered Ridge,
Hub, Vale or another way, even the first calculation is not reproducible.
If the owner map is missing, the field value can still be correct while
the capacity account remains undefined. An honest report can say both
things at once rather than treating the missing layer as permission to
guess.

\section{A final diagnostic exercise}
Imagine a report for the three-child source example gives the group
value thirteen, says every owner starts at zero, and claims that a
closing receipt proves fair compensation. Identify three separate issues.
First, thirteen is the sum of same-baseline singletons; the actual joint
field value is ten. Second, balances alone do not identify owners or the
rule mapping path values to accounts. Third, no corrected sum would prove
fairness. Those defects require different repairs.

Now imagine the report instead gives ten and correct path values but
executes only two of the three children. The formula has not failed.
The event identity has: a different increment was performed. A receipt
must correspond to that actual authorized event or the system must refuse
under its declared deviation procedure. Carry this distinction into the
software chapters. Their central job is to prevent a correct calculation
from becoming a misleading record of a different occurrence.

\section{A route-and-repair reading: the remote pump seal}
The original laboratory also used a remote pump seal to explain genuine
need, route choice and live-state replanning. That lesson deserves to
survive the change in potential. A water-treatment plant needs a working
seal before a deadline; rail and air offer different delays and physical
burdens. A route is not valuable or worthless merely because it is long.
The represented endpoint, timing, feasibility and actual process matter.

This repair is not an action inside the first homogeneous-scalar Gaussian
world. A seal is a discrete functional item, not divisible water, and a
plant's working condition is not automatically a scalar stock coordinate.
A broader domain must declare the transformation and its functional
representation. Keeping that boundary visible preserves the example's
meaning without pretending the first experiment already models maintenance.

The original example assigns illustrative planning burdens of nine to
rail and seventeen to air before additional packaging. It also records
usable factors $0.99\cdot0.99=0.9801$ and
$0.99\cdot0.98=0.9702$. A fraction of a seal may be useless, so
these products cannot simply be treated as divisible functional delivery.
Packaging, reliability and item-count assumptions must be explicit.
The historical exercise later assumes successful functional delivery;
that is a teaching assumption, not a measured guarantee that probability
one can be certified in a real transport system.

Suppose a flood closes the planned rail connection after the first segment.
The planned endpoint is no longer the live physical situation. Replanning
must start from the actual hub state and the new route constraints.
An old positive forecast is not a receipt for an unperformed route. If
the flood also changes represented physical state or parameters, that
external change needs its own account; it cannot be hidden inside an
actor's next value.

The original's separate fixed-state accounting exercise lists potentials
$80,76,74,68,5$ and burdens $3,1,8,4$. Under one fixed potential
and with every intervening change included in those declared actions,
the signed lines are $1,1,-2,59$, summing to fifty-nine. Endpoint
closure gives the same $80-5-16=59$. If an unrecorded external change
occurs between those states, this simple actor-only telescoping would
need the external correction derived earlier. The arithmetic cannot
certify that all such events were recorded.

Repeated repairs can motivate separate studies of spare inventory,
local manufacture, standardization or predictive maintenance. They do
not justify crediting today's actor with all forecast future savings.
Those interventions have their own physical costs, uncertainty and
institutional questions. The future Part VI develops routes and
transformations; Part VII develops explicit coordination policies.
Neither is hidden inside the present random actor.

\section{Proof workshop: rebuild the argument in eight passes}
The goal of this workshop is not to memorize eight displayed equations.
It is to learn how to reconstruct a result and state where it stops.
For each pass, first write the objects and assumptions, then the algebra,
then one conclusion that the algebra does not establish.

\subsection{Pass 1: conservation from the transformation matrix}
Begin with $x^+=x+Sq$. Multiply by the row of ones. If every column
of $S$ sums to zero, the total before and after is equal for every
quantity vector. This is a structural statement about the declared
transformation matrix. To complete a physical validation, separately
check nonnegative coordinates, allowed quantities and any joint constraint.
An impossible negative source is not repaired by the conservation proof.

Now replace one destination coefficient by $\eta<1$. The old
column sum changes. State explicitly whether the missing carrier enters
a loss coordinate or crosses an external boundary. Do not repair the
mass equation by adding a burden number to it: burden and physical
stock have different units and meanings.

\subsection{Pass 2: differentiate the declared potential}
Start with one term $\tfrac12(x_i-x_i^*)^2/\sigma_i^2$.
Keep the reference and scale fixed while differentiating. Recover
$\mu_i=(x_i-x_i^*)/\sigma_i^2$ and curvature
$1/\sigma_i^2$. Explain why differentiating a moving scale or
reference would create additional terms. The fixed-parameter condition
is part of the calculation, not a footnote to discard afterward.

Check the units: marginal has inverse-stock units and curvature has
inverse-stock-squared units under dimensionless normalization. Finally
explain why a derivative does not identify an actor's preferred action.
It describes a directional property of a function before a choice rule
has been defined.

\subsection{Pass 3: expand a finite movement}
Write $y=z-x^*$ and expand
$\tfrac12(y+\delta)^TH(y+\delta)$. The two cross terms
are equal because $H$ is symmetric. Subtract the initial potential
and change the sign to obtain exact field benefit. The positive
quadratic term in potential becomes a negative curvature correction
in benefit. Check the result at zero increment and at the reference.

At the reference a nonzero net increment has strictly negative field
value. A nonempty cancelling group has zero net increment instead.
The proof distinguishes those cases and therefore cannot be paraphrased
as ``every nonempty group at equilibrium has negative value.''

\subsection{Pass 4: cancel unchanged factors}
Partition the potential's factors into those touched by the action's
write support and those untouched. Subtract before and after. Every
untouched factor cancels exactly, leaving the complete affected-factor
sum. For a factor involving several cells, include all arguments needed
to evaluate it even when some are not changed.

The proof establishes the locality of this value calculation. It does
not prove that a world-wide group enumerator is decentralized. Write
that limitation beside the theorem so the word local does not quietly
change meaning between a mathematical and an architectural claim.

\subsection{Pass 5: recover the path integral}
Let $\gamma(s)=z+s\delta$. Apply the chain rule to
$V(\gamma(s))$ and integrate from zero to one. The endpoint
difference follows from the fundamental theorem of calculus. In the
quadratic case the gradient is affine, so integrating it gives the
mean of its endpoint gradients exactly.

Then state which path meaning is in force. Is it a constant-rate
model-realized trajectory or merely an endpoint interpolation? The
integral identity does not choose. An observed real-world path requires
more than an equality between two endpoint calculations.

\subsection{Pass 6: compose children before dividing value}
Expand the square of $\sum_a\delta_a$. Keep all pair cross
terms. Their omission is the precise defect in summing independent
same-baseline singletons. Next integrate each child against the same
group gradient. Summing those child integrals recovers the integral
of the net increment and hence group closure.

The closure proof has no requester, supplier, beneficiary or owner in
it. That absence is informative. It is why choosing an owner map is
another modelling step rather than a missing algebraic manipulation.
Adding a label to a child does not add a theorem of entitlement.

\subsection{Pass 7: telescope a sequence, then insert forcing}
Write three consecutive field differences and cancel intermediate
potentials visibly. Extend the pattern to any finite sequence. For
a closed cycle the field sum vanishes; nonnegative separately declared
burdens make net value nonpositive. Do not replace the live baseline
of every line by the first state.

Now insert a forcing step before each actor epoch. Add and subtract
the pre-forcing potential. The new external terms explain why an
actor-only sum need not equal the overall endpoint difference. This
is not a failure of telescoping; it is the account of an additional
kind of change that must not be credited as actor work.

\subsection{Pass 8: derive the account invariant conditionally}
Assume an explicitly adopted owner rule whose signed changes sum to
the field value, and a signed external ledger updated only by forcing.
Add the potential change, balance change and negative audit change.
All terms cancel under those definitions. Record the invariant and
its initial constant.

Finally construct the restoring-and-reversing pair already shown in
the laboratory. Both events satisfy the invariant, yet the second
undoes recovery by spending prior capacity. This last step prevents
the workshop from ending with an unjustified dynamical conclusion.
The invariant constrains accounting; the declared stochastic process
still determines which permitted events occur.

\section{Ten mistakes a laboratory should expose}
\textbf{1. Omitting efficiency outside the lossless domain.}
The destination increment changes, so both the physical balance and
the directional derivative change. A lossless example cannot certify
a loss-aware implementation.

\textbf{2. Valuing from the pre-forcing state.}
This folds nature's change into the actor's line. Keep the external
record and the post-forcing baseline separate.

\textbf{3. Ranking by value per unit without declaring an optimizer.}
A ratio may be a legitimate diagnostic or policy in another model.
Using it to select actions changes the current unweighted-random
question, even if the finite quote itself is exact.

\textbf{4. Settling the requested quantity instead of the accepted event.}
The same action name can conceal a different increment. Bind quantity
and baseline to the event identity, not just to a display label.

\textbf{5. Counting one receipt twice.}
Arithmetic correctness of each copy does not make the second copy a
second physical occurrence. An account needs exactly-once identity and
replay rules as well as a correct subtraction.

\textbf{6. Pricing shared-source children independently.}
The common source changes only once in the joint endpoint. Restore
the cross terms or evaluate the complete endpoint directly.

\textbf{7. Giving the full source budget to every edge.}
Individually valid requests can be jointly impossible. A physical
budget has one scope; duplicating its record does not duplicate stock.

\textbf{8. Reading continuous dissipation as a finite-step guarantee.}
The exact curvature term survives discretization. Check the accepted
finite increment and physical feasibility independently.

\textbf{9. Projecting an invalid endpoint while keeping its quote.}
A clipped successor is another physical map. It needs the value of
that map and an explicit permission rule, not a repaired-looking state
attached to the unrepaired record.

\textbf{10. Inferring forever from a finite study.}
A completed horizon establishes observations on that registered horizon.
Unlimited-time claims need their own theorem or evidential argument.
Neither a large test count nor a smooth plot removes that boundary.

These mistakes are not ten ways the same equation fails. They are
ten ways a reader or implementation can quietly substitute a different
object for the one the equation describes. The following software
chapters turn that diagnosis into inspectable responsibilities.
